\documentclass[10pt,a4paper]{amsart}
\usepackage[margin=19mm]{geometry}
\usepackage{amsmath,amssymb,mathtools}
\usepackage[scr=rsfs,cal=euler]{mathalfa}
\usepackage{xfrac}
\usepackage[unicode,hidelinks,bookmarks=false]{hyperref}
\newcommand{\ie}{i.e.\ }
\newcommand{\cf}{cf.\ }
\newcommand{\resp}{respectively\ }
\newcommand{\Subsection}{\subsection}
\providecommand{\nobreakdash}{\nobreak\hbox{-}}
\setlength{\emergencystretch}{2em}
\multlinegap0pt
\usepackage{xcolor}
\usepackage{amssymb}
\usepackage{mathtools}
\usepackage{bm}
\usepackage[shortlabels]{enumitem}
\newtheorem{assumption}{Assumption}
\newtheorem{theorem}{Theorem}
\newtheorem{lemma}{Lemma}
\title{Continuous-Time Analysis for Minimax and Bilevel Problems}
\author{Hyunwoo Lee, Jeongyeol Kwon, Dohyun Kwon}
\date{Source: arXiv:2605.20898v1 (CC BY 4.0)}
\begin{document}
\maketitle

\noindent\emph{Source and licence.} Continuous-Time Analysis for Minimax and Bilevel Problems. Version arXiv:2605.20898v1 (CC BY 4.0), distributed by arXiv under the Creative Commons Attribution 4.0 International licence, http://creativecommons.org/licenses/by/4.0/, as recorded in the arXiv licence metadata for this exact version and shown on the abstract page https://arxiv.org/abs/2605.20898v1. Full licence notice: \texttt{licenses/arxiv-2605.20898-CC-BY-4.0.txt}.\par\smallskip

\noindent\textit{Editorial scope.} Counted unit: the article's lemma lem:bilevel-residuals (source0.tex lines 2130--2163). The article's complete original proof of it is delivered with it (source0.tex lines 2165--2242, one \texttt{proof} environment of 78 lines). No mathematics is written, repaired or re-derived: the statement and the proof are the article's own lines, in the article's own notation, and no retained line is substituted or re-flowed.  Retained uncounted as the source-local context the proof needs: lines 165--171; lines 298--308; lines 382--388; lines 663--671; lines 685--695.  Nothing was omitted from any retained span, so the package carries no omission record.  No retained line contains a citation, so the wrapper carries no bibliography and no bibliography entry.  No retained line contains a figure environment, an image-inclusion command or drawing code, so no image is delivered and the package declares no asset dependency.  Editorial formatting only, in addition: an AMS \texttt{amsart} preamble that restates the article's own declarations and macros used by the retained lines, namely the environments \texttt{assumption}, \texttt{theorem}, \texttt{lemma} on separate counters, together with the standard packages those lines need.  The retained environments print in this document as Assumption 1, Theorem 1, Lemma 1 in order of appearance, whereas the article prints the same items with the numbering of its own full document.  The complete native source of the article as distributed in the arXiv e-print for this exact version is bundled unchanged as \texttt{sources/201029/source0.tex}.
\begin{equation}\label{eq:problem2}
\tag{\textbf{P2}}
%\begin{aligned}
\min_{x\in X}\quad  F(x)\ :=\ f\bigl(x,y^{*}(x)\bigr) 
\quad \text{ s.t. }\quad  y^{*}(x)\ \in\ \arg\min_{y\in Y}\ g(x,y),
%\end{aligned}
\end{equation}

\begin{assumption}\label{ass:2}
We consider \ref{eq:problem2} with continuously differentiable functions $f$ and $g$.
There exist constants $\mu>0$, $\rho\ge 0$, and $L_{fx},L_{gx}\ge 0$ such that:
\begin{enumerate}%[leftmargin=*]
  \item For every $x$, $g(x,\cdot)$ is $\mu$-strongly convex.
  \item $\nabla_x f(x,y)$ is $L_{fx}$-Lipschitz in $y$, uniformly over $x$.
  \item $\nabla_x g(x,y)$ is $L_{gx}$-Lipschitz in $y$, uniformly over $x$.
  \item $\mathcal{L}_{\lambda,\inf}^{*}:=\inf_{x\in\mathbb{R}^{d_x}} \mathcal L_\lambda^{*}(x) > -\infty.$
  \item For every $x$, $f(x,\cdot)$ is $\rho$--weakly convex.
\end{enumerate}
\end{assumption}

\begin{equation}\label{eq:flow-bilevel-unified}
\begin{aligned}
\dot x = -\nabla_x E_\lambda(x,y,z),\
\dot y = -\delta\,\nabla_y E_\lambda(x,y,z),\
\dot z = \eta\,\nabla_z E_\lambda(x,y,z).
\end{aligned}
\end{equation}

\begin{align}
\delta_0(\lambda)
\ &:=\ \frac{C_y(\lambda)^2}{\alpha}
\left(\frac{1-\alpha}{2}+\frac{\theta|\alpha-\beta|}{2}\right), \
\eta_0
\ :=\ \frac{C_z^2}{\beta}
\left(\beta+\frac{(1-\beta)^2}{1-\alpha}+\frac{|\alpha-\beta|}{2\theta}\right).
\label{eq:eta0-bilevel}
\end{align}

\begin{theorem}\label{thm:fixed-lambda-simple}
Assume Assumption~\ref{ass:2} holds. Fix $\lambda>\rho/\mu$, $\alpha,\beta\in(0,1)$, and $\theta>0$.
Consider any solution $(x(t),y(t),z(t))$ of \eqref{eq:flow-bilevel-unified}.
If the time-scale parameters satisfy $\delta \ \ge\ \delta_0(\lambda)$ and $\eta \ \ge\ \eta_0$, where $\delta_0(\lambda)$ and $\eta_0$ are defined in \eqref{eq:eta0-bilevel},
then for $t\sim\mathrm{Unif}[0,T]$,
$
\mathbb E_t\Big[\bigl\|\nabla_x \mathcal L_\lambda^{*}(x(t))\bigr\|^{2}\Big]
\ \le\ \frac{W(x(0),y(0),z(0))-\mathcal L_{\lambda,\inf}^{*}}{c_1\,T},
$
where $c_1= (1-\alpha)/4$ .
\end{theorem}

\begin{lemma}\label{lem:bilevel-residuals}
Assume the conditions of Theorem~\ref{thm:fixed-lambda-simple}.
Recall the residuals along the trajectory $(x(t),y(t),z(t))$ of \eqref{eq:flow-bilevel-unified}:
\[
r_y(t):=\nabla_y \mathcal L_\lambda(x(t),y(t))
=\nabla_y f(x(t),y(t))+\lambda\nabla_y g(x(t),y(t)),
\qquad
r_z(t):=\nabla_y g(x(t),z(t)).
\]
Let $\delta_0(\lambda)$ and $\eta_0$ be defined in \eqref{eq:eta0-bilevel}.
Suppose that
\[
\delta>\delta_0(\lambda),
\qquad
\eta>\eta_0,
\]
and define the positive margins
\[
c_y(\delta,\lambda):=\alpha\bigl(\delta-\delta_0(\lambda)\bigr),
\qquad
c_z(\eta,\lambda):=\beta\lambda^2\bigl(\eta-\eta_0\bigr).
\]
Then for $t\sim\mathrm{Unif}[0,T]$,
\begin{align}
\mathbb E_t\bigl[\|r_y(t)\|^2\bigr]
&\le
\frac{W(x(0),y(0),z(0))-\mathcal L_{\lambda,\inf}^{*}}{c_y(\delta,\lambda)\,T},
\label{eq:ergodic-F-lemma}\\
\mathbb E_t\bigl[\|r_z(t)\|^2\bigr]
&\le
\frac{W(x(0),y(0),z(0))-\mathcal L_{\lambda,\inf}^{*}}{c_z(\eta,\lambda)\,T}.
\label{eq:ergodic-Gz-lemma}
\end{align}
\end{lemma}

\begin{proof} 
The proof of Theorem~\ref{thm:fixed-lambda-simple} establishes the inequality: there exist constants $c_1>0$, $c_2>0$, and $c_3>0$ such that for a.e.\ $t\ge 0$,
\begin{equation}\label{eq:lyap-refined}
\frac{d}{dt}W(x(t),y(t),z(t))
\le
-c_1\|\nabla_x \mathcal L_\lambda^*(x(t))\|^2
-c_2\|r_y(t)\|^2
-c_3\|r_z(t)\|^2.
\end{equation}
Moreover, with the explicit definitions in the proof,
\[
c_2=\alpha\delta-K_yC_y(\lambda)^2,
\qquad
c_3=\beta\eta\lambda^2-K_zC_z^2,
\]
and the thresholds $\delta_0(\lambda)=K_yC_y(\lambda)^2/\alpha$ and $\eta_0=K_zC_z^2/(\beta\lambda^2)$
are chosen precisely so that $c_2>0$ and $c_3>0$. Under the strict inequalities $\delta>\delta_0(\lambda)$ and $\eta>\eta_0$, we can rewrite
\[
c_2=\alpha\delta-\alpha\delta_0(\lambda)=\alpha(\delta-\delta_0(\lambda))=:c_y(\delta,\lambda)>0,
\]
\[
c_3=\beta\eta\lambda^2-\beta\eta_0\lambda^2=\beta\lambda^2(\eta-\eta_0)=:c_z(\eta,\lambda)>0.
\]
Substituting these into \eqref{eq:lyap-refined} yields
\begin{equation}\label{eq:lyap-refined-margins}
\frac{d}{dt}W(x(t),y(t),z(t))
\le
-c_1\|\nabla_x \mathcal L_\lambda^*(x(t))\|^2
-c_y(\delta,\lambda)\|r_y(t)\|^2
-c_z(\eta,\lambda)\|r_z(t)\|^2.
\end{equation}

Integrating \eqref{eq:lyap-refined-margins} over $[0,T]$ gives
\begin{align*}
W(x(T),y(T),z(T)) - W(x(0),y(0),z(0))
&\le
- c_1\int_0^T \|\nabla_x \mathcal L_\lambda^*(x(t))\|^2\,dt \\
&\quad
- c_y(\delta,\lambda)\int_0^T \|r_y(t)\|^2\,dt
- c_z(\eta,\lambda)\int_0^T \|r_z(t)\|^2\,dt.
\end{align*}
Dropping the nonpositive term involving $\|\nabla_x \mathcal L_\lambda^*(x(t))\|^2$ yields the simpler inequality
\begin{equation}\label{eq:integral-two-residuals}
c_y(\delta,\lambda)\int_0^T \|r_y(t)\|^2\,dt
+
c_z(\eta,\lambda)\int_0^T \|r_z(t)\|^2\,dt
\le
W(x(0),y(0),z(0)) - W(x(T),y(T),z(T)).
\end{equation}
By definition,
\[
W(x,y,z)
=
\mathcal L_\lambda^*(x)+\alpha\widetilde H_\lambda(x,y)+\beta\lambda H(x,z),
\]
and $\widetilde H_\lambda\ge 0$, $H\ge 0$. Hence $W(x,y,z)\ge \mathcal L_\lambda^*(x)$ for all $(x,y,z)$.
Since $\mathcal L_{\lambda,\inf}^*:=\inf_x \mathcal L_\lambda^*(x)>-\infty$ by Assumption~\ref{ass:2}.4, we have
\[
W(x(T),y(T),z(T))\ge \mathcal L_\lambda^*(x(T))\ge \mathcal L_{\lambda,\inf}^*.
\]
Combining with \eqref{eq:integral-two-residuals} yields the two separate bounds
\[
c_y(\delta,\lambda)\int_0^T \|r_y(t)\|^2\,dt
\le
W(x(0),y(0),z(0))-\mathcal L_{\lambda,\inf}^*,
\qquad
c_z(\eta,\lambda)\int_0^T \|r_z(t)\|^2\,dt
\le
W(x(0),y(0),z(0))-\mathcal L_{\lambda,\inf}^*.
\]
If $t\sim\mathrm{Unif}[0,T]$, then
\[
\mathbb E_t[\|r_y(t)\|^2]=\frac{1}{T}\int_0^T \|r_y(s)\|^2\,ds,
\qquad
\mathbb E_t[\|r_z(t)\|^2]=\frac{1}{T}\int_0^T \|r_z(s)\|^2\,ds.
\]
Dividing the last two inequalities by $T$ gives \eqref{eq:ergodic-F-lemma}--\eqref{eq:ergodic-Gz-lemma}.
\end{proof}

\end{document}
